\documentclass[10pt,letterpaper]{article}
\usepackage[top=0.85in,left=2.75in,footskip=0.75in]{geometry}
% amsmath and amssymb packages, useful for mathematical formulas and symbols
\usepackage{amsmath,amssymb}
% Use adjustwidth environment to exceed column width (see example table in text)
\usepackage{changepage}
% Use Unicode characters when possible
\usepackage[utf8x]{inputenc}
% textcomp package and marvosym package for additional characters
\usepackage{textcomp,marvosym}
% cite package, to clean up citations in the main text. Do not remove.
\usepackage{cite}
% Use nameref to cite supporting information files (see Supporting Information section for more info)
\usepackage{nameref,hyperref}
% line numbers
\usepackage[right]{lineno}
% ligatures disabled
\usepackage{microtype}
\DisableLigatures[f]{encoding = *, family = * }
% color can be used to apply background shading to table cells only
\usepackage[table]{xcolor}
% array package and thick rules for tables
\usepackage{array}
% create "+" rule type for thick vertical lines
\newcolumntype{+}{!{\vrule width 2pt}}
% create \thickcline for thick horizontal lines of variable length
\newlength\savedwidth
\newcommand\thickcline[1]{%
  \noalign{\global\savedwidth\arrayrulewidth\global\arrayrulewidth 2pt}%
  \cline{#1}%
  \noalign{\vskip\arrayrulewidth}%
  \noalign{\global\arrayrulewidth\savedwidth}%
}
% \thickhline command for thick horizontal lines that span the table
\newcommand\thickhline{\noalign{\global\savedwidth\arrayrulewidth\global\arrayrulewidth 2pt}%
\hline
\noalign{\global\arrayrulewidth\savedwidth}}

% Remove comment for double spacing
%\usepackage{setspace} 
%\doublespacing

% Text layout
\raggedright
\setlength{\parindent}{0.5cm}
\textwidth 5.25in 
\textheight 8.75in

% Bold the 'Figure #' in the caption and separate it from the title/caption with a period
% Captions will be left justified
\usepackage[aboveskip=1pt,labelfont=bf,labelsep=period,justification=raggedright,singlelinecheck=off]{caption}
\renewcommand{\figurename}{Fig}

% Use the PLoS provided BiBTeX style
\bibliographystyle{plos2015}

% Remove brackets from numbering in List of References
\makeatletter
\renewcommand{\@biblabel}[1]{\quad#1.}
\makeatother

% Leave date blank
\date{}

% Header and Footer with logo
\usepackage{lastpage,fancyhdr,graphicx}
\usepackage{epstopdf}
\pagestyle{myheadings}
\pagestyle{fancy}
\fancyhf{}
\setlength{\headheight}{27.023pt}
\lhead{\includegraphics[width=2.0in]{PLOS-submission.eps}}
\rfoot{\thepage/\pageref{LastPage}}
\renewcommand{\footrule}{\hrule height 2pt \vspace{2mm}}
\fancyheadoffset[L]{2.25in}
\fancyfootoffset[L]{2.25in}
\lfoot{\sf PLOS}

%% Include all macros below

\newcommand{\lorem}{{\bf LOREM}}
\newcommand{\ipsum}{{\bf IPSUM}}

% added by sgy
\usepackage{url}
\newcommand{\col}[2][red]{\textcolor{#1}{#2}}
\newcommand{\bi}{\begin{itemize}}
\newcommand{\ei}{\end{itemize}}
\newcommand{\package}{\emph{ITHIM}}
\newcommand{\ithim}{\emph{ITHIM}}
\newcommand{\R}{\emph{R}}
\newcommand{\sgy}[1]{{\itshape\col{#1}}}
\newcommand{\XX}{\sgy{XX}}
\newcommand{\webpage}{\url{http://www.ithim.ghi.wisc.edu}}
\newcommand{\af}{\mathrm{PAF}}
\newcommand{\mets}{MET-hrs./week}
\newcommand{\pEm}[1]{$\mathrm{PM}_{#1}$}
\newcommand{\vmt}{$\mathrm{VMT}$}
\newcommand{\logNormal}{log-normal}


%% END MACROS SECTION

\begin{document}
\vspace*{0.2in}

% Title must be 250 characters or less.
\begin{flushleft}
{\Large
\textbf\newline{Improved Implementation of Comparative Risk
  Assessment for Health Benefits of Active Transportation Using
  National Datasets} % Please use "title case" (capitalize all terms in the title except conjunctions, prepositions, and articles).
}
\newline
% Insert author names, affiliations and corresponding author email (do not include titles, positions, or degrees).
\\
Samuel Younkin\textsuperscript{1},\ %,2\Yinyang},
Jason Vargo\textsuperscript{1},\ %\Yinyang},
Jonathan Patz\textsuperscript{1}%,3\textcurrency},
%Name4 Surname\textsuperscript{2},
%#Name5 Surname\textsuperscript{2\ddag},
%Name6 Surname\textsuperscript{2\ddag},
%Name7 Surname\textsuperscript{1,2,3*},
%with the Global Health Institute\textsuperscript{\textpilcrow}
\\
\bigskip
\textbf{1} Global Health Institute, University of Wisconsin{\textendash}Madison, Madison, WI, USA
\\
%\textbf{2} Affiliation Dept/Program/Center, Institution Name, City, State, Country
%\\
%\textbf{3} Affiliation Dept/Program/Center, Institution Name, City, State, Country
%\\
\bigskip

% Insert additional author notes using the symbols described below. Insert symbol callouts after author names as necessary.
% 
% Remove or comment out the author notes below if they aren't used.
%
% Primary Equal Contribution Note
%\Yinyang These authors contributed equally to this work.

% Additional Equal Contribution Note
% Also use this double-dagger symbol for special authorship notes, such as senior authorship.
% \ddag These authors also contributed equally to this work.

% Current address notes
% \textcurrency Current Address: GHI UW{\textendash}Madison, 1300
% University Ave., Madison, WI, 53706, USA
% WI, USA % change symbol to "\textcurrency a" if more than one current address note
% \textcurrency b Insert second current address 
% \textcurrency c Insert third current address

% Deceased author note
%\dag Deceased

% Group/Consortium Author Note
%\textpilcrow Membership list can be found in the Acknowledgments section.

% Use the asterisk to denote corresponding authorship and provide email address in note below.
%* correspondingauthor@institute.edu

% Please keep the abstract below 300 words
% Use "Eq" instead of "Equation" for equation citations.

% Results and Discussion can be combined.

\end{flushleft}

\section*{Abstract}
\section{Summary:}

\section{Contact:}
\href{javargo@wisc.edu}{javargo@wisc.edu}



% Please keep the Author Summary between 150 and 200 words
% Use first person. PLOS ONE authors please skip this step. 
% Author Summary not valid for PLOS ONE submissions.   
% \section*{Author Summary}
% Lorem ipsum dolor sit amet, consectetur adipiscing elit. Curabitur
eget porta erat. Morbi consectetur est vel gravida
pretium. Suspendisse ut dui eu ante cursus gravida non sed sem. Nullam
sapien tellus, commodo id velit id, eleifend volutpat quam. Phasellus
mauris velit, dapibus finibus elementum vel, pulvinar non tellus. Nunc
pellentesque pretium diam, quis maximus dolor faucibus id. Nunc
convallis sodales ante, ut ullamcorper est egestas vitae. Nam sit amet
enim ultrices, ultrices elit pulvinar, volutpat risus.


\linenumbers

\section*{Introduction}
Hello world.


\section*{Materials and Methods}
The ITHIM model uses the comparative risk assesment framework to
estimate the population attributable fraction, $\af$.  In our case the
$\af$ represents the proportional change in national disease burden if
age and sex-specific active transport time means were changed.

\begin{equation}
\af = \frac{\int \! R(x)P(x) \, \mathrm{d}x  - \int \! R(x)Q(x) \, \mathrm{d}x}{\int \! R(x)P(x) \, \mathrm{d}x}
\end{equation}

Here $R(x)$ represents the risk when the individual has exposure $x$,
physical activity (\mets).  $P(x)$ and $Q(x)$ denote the population
density of individuals with exposure $x$, in the baseline and
alternate scenarios, respectively.  The exposure variable, $x$, is
modeled as the sum of two independent random variables, travel-related
exposure and non-travel-related exposure.  Non-travel-related exposure
is modeled with a \logNormal{} distribution with the age-sex specific
mean ratios for non-travel-related exposure are estimated from the
ATUS \cite{ATUS}.  The referent group mean is treated as a variable
and the coefficient of variation is treated as constant across age-sex
strata and estimated with ATUS.  We see in Figure \ref{dalyFigure}
that varying the mean values for non-travel-related exposure does not
greatly affect the estimate for disease burden.  The distribution for
travel-related exposure is estimated from the time spent walking or
cycling, i.e., active transport time, which as in the original ITHIM
model has a \logNormal{} distribution with constant coefficient of
variation across age-sex strata.  Travel-related exposure is estimated
from active transport time using assumptions about how much physical
activity is required for cycling and walking.


\section*{Results}
\begin{figure}[t]
    \centerline{\includegraphics[width=\textwidth]{./figures/fig2}}
    \caption{}\label{dalyFigure}
\end{figure}

\begin{figure}[t]
  \centerline{\includegraphics[width=\textwidth]{./figures/fig3.pdf}}
    \caption{Mean travel time by mode (walking/cycling) for each of
      seven states.  National means are displayed on the
      right.}\label{meanMatrices}
\end{figure}

\begin{figure}[t]
  \centerline{\includegraphics[width=\textwidth]{./figures/fig4.pdf}}
    \caption{}\label{California}
\end{figure}

We used the National Household Transportation Survey to estimate the
mean walk and cycle times for individuals that live in metropolitan
areas for each of seven states (CA, TX, NY, FL, VA, NC and AZ)
\cite{NHTS}.  These states were selected becaused they had sufficient
sample size for reasonable estimates of age-sex specific
travel-related walking and cycling means, Figure \ref{meanMatrices}.

Of these seven states one was found to show active transport behavior
that, if adopted nation-wide, would avert 75-100,000 DALYs overall.
Values for national disease-specific DALYs were found using the CDC
Wonder database \cite{CDCWonder}.  California stood out among the
estimates for national DALYs averted.  Most of the decrease in DALYs
is due to prevention of dementia in the oldest age group, $80+$ years
old.  California exhibits much greater walking and cycling means in
this age group than the national average. In particular we see the
cycling mean We also see prevention of depression and diabetes in the
middle age groups due to increased active transport time.


\section*{Discussion}
Nulla mi mi, venenatis sed ipsum varius, Table~\ref{table1} volutpat
euismod diam. Proin rutrum vel massa non gravida. Quisque tempor sem
et dignissim rutrum. Lorem ipsum dolor sit amet, consectetur
adipiscing elit. Morbi at justo vitae nulla elementum commodo eu id
massa. In vitae diam ac augue semper tincidunt eu ut eros. Fusce
fringilla erat porttitor lectus cursus, vel sagittis arcu
lobortis. Aliquam in enim semper, aliquam massa id, cursus
neque. Praesent faucibus semper libero~\cite{bib3}.


\section*{Conclusion}
CO\textsubscript{2} Maecenas convallis mauris sit amet sem ultrices gravida. Etiam eget sapien nibh. Sed ac ipsum eget enim egestas ullamcorper nec euismod ligula. Curabitur fringilla pulvinar lectus consectetur pellentesque. Quisque augue sem, tincidunt sit amet feugiat eget, ullamcorper sed velit. 

Sed non aliquet felis. Lorem ipsum dolor sit amet, consectetur adipiscing elit. Mauris commodo justo ac dui pretium imperdiet. Sed suscipit iaculis mi at feugiat. Ut neque ipsum, luctus id lacus ut, laoreet scelerisque urna. Phasellus venenatis, tortor nec vestibulum mattis, massa tortor interdum felis, nec pellentesque metus tortor nec nisl. Ut ornare mauris tellus, vel dapibus arcu suscipit sed. Nam condimentum sem eget mollis euismod. Nullam dui urna, gravida venenatis dui et, tincidunt sodales ex. Nunc est dui, sodales sed mauris nec, auctor sagittis leo. Aliquam tincidunt, ex in facilisis elementum, libero lectus luctus est, non vulputate nisl augue at dolor. For more information, see \nameref{S1_Appendix}.


\section*{Supporting Information}

CO\textsubscript{2} Maecenas convallis mauris sit amet sem ultrices gravida. Etiam eget sapien nibh. Sed ac ipsum eget enim egestas ullamcorper nec euismod ligula. Curabitur fringilla pulvinar lectus consectetur pellentesque. Quisque augue sem, tincidunt sit amet feugiat eget, ullamcorper sed velit. 

Sed non aliquet felis. Lorem ipsum dolor sit amet, consectetur adipiscing elit. Mauris commodo justo ac dui pretium imperdiet. Sed suscipit iaculis mi at feugiat. Ut neque ipsum, luctus id lacus ut, laoreet scelerisque urna. Phasellus venenatis, tortor nec vestibulum mattis, massa tortor interdum felis, nec pellentesque metus tortor nec nisl. Ut ornare mauris tellus, vel dapibus arcu suscipit sed. Nam condimentum sem eget mollis euismod. Nullam dui urna, gravida venenatis dui et, tincidunt sodales ex. Nunc est dui, sodales sed mauris nec, auctor sagittis leo. Aliquam tincidunt, ex in facilisis elementum, libero lectus luctus est, non vulputate nisl augue at dolor. For more information, see \nameref{S1_Appendix}.


\section*{Acknowledgments}

\nolinenumbers

\bibliography{../tex/ITHIM}

\end{document}


% Template for PLoS
% Version 3.3 June 2016
%
% % % % % % % % % % % % % % % % % % % % % %
%
% -- IMPORTANT NOTE
%
% This template contains comments intended 
% to minimize problems and delays during our production 
% process. Please follow the template instructions
% whenever possible.
%
% % % % % % % % % % % % % % % % % % % % % % % 
%
% Once your paper is accepted for publication, 
% PLEASE REMOVE ALL TRACKED CHANGES in this file 
% and leave only the final text of your manuscript. 
% PLOS recommends the use of latexdiff to track changes during review, as this will help to maintain a clean tex file.
% Visit https://www.ctan.org/pkg/latexdiff?lang=en for info or contact us at latex@plos.org.
%
%
% There are no restrictions on package use within the LaTeX files except that 
% no packages listed in the template may be deleted.
%
% Please do not include colors or graphics in the text.
%
% The manuscript LaTeX source should be contained within a single file (do not use \input, \externaldocument, or similar commands).
%
% % % % % % % % % % % % % % % % % % % % % % %
%
% -- FIGURES AND TABLES
%
% Please include tables/figure captions directly after the paragraph where they are first cited in the text.
%
% DO NOT INCLUDE GRAPHICS IN YOUR MANUSCRIPT
% - Figures should be uploaded separately from your manuscript file. 
% - Figures generated using LaTeX should be extracted and removed from the PDF before submission. 
% - Figures containing multiple panels/subfigures must be combined into one image file before submission.
% For figure citations, please use "Fig" instead of "Figure".
% See http://journals.plos.org/plosone/s/figures for PLOS figure guidelines.
%
% Tables should be cell-based and may not contain:
% - spacing/line breaks within cells to alter layout or alignment
% - do not nest tabular environments (no tabular environments within tabular environments)
% - no graphics or colored text (cell background color/shading OK)
% See http://journals.plos.org/plosone/s/tables for table guidelines.
%
% For tables that exceed the width of the text column, use the adjustwidth environment as illustrated in the example table in text below.
%
% % % % % % % % % % % % % % % % % % % % % % % %
%
% -- EQUATIONS, MATH SYMBOLS, SUBSCRIPTS, AND SUPERSCRIPTS
%
% IMPORTANT
% Below are a few tips to help format your equations and other special characters according to our specifications. For more tips to help reduce the possibility of formatting errors during conversion, please see our LaTeX guidelines at http://journals.plos.org/plosone/s/latex
%
% For inline equations, please be sure to include all portions of an equation in the math environment.  For example, x$^2$ is incorrect; this should be formatted as $x^2$ (or $\mathrm{x}^2$ if the romanized font is desired).
%
% Do not include text that is not math in the math environment. For example, CO2 should be written as CO\textsubscript{2} instead of CO$_2$.
%
% Please add line breaks to long display equations when possible in order to fit size of the column. 
%
% For inline equations, please do not include punctuation (commas, etc) within the math environment unless this is part of the equation.
%
% When adding superscript or subscripts outside of brackets/braces, please group using {}.  For example, change "[U(D,E,\gamma)]^2" to "{[U(D,E,\gamma)]}^2". 
%
% Do not use \cal for caligraphic font.  Instead, use \mathcal{}
%
% % % % % % % % % % % % % % % % % % % % % % % % 
%
% Please contact latex@plos.org with any questions.
%
% % % % % % % % % % % % % % % % % % % % % % % %

% \subsection*{Etiam eget sapien nibh.}

% % For figure citations, please use "Fig" instead of "Figure".
% Nulla mi mi, Fig~\ref{fig1} venenatis sed ipsum varius, volutpat euismod diam. Proin rutrum vel massa non gravida. Quisque tempor sem et dignissim rutrum. Lorem ipsum dolor sit amet, consectetur adipiscing elit. Morbi at justo vitae nulla elementum commodo eu id massa. In vitae diam ac augue semper tincidunt eu ut eros. Fusce fringilla erat porttitor lectus cursus, \nameref{S1_Video} vel sagittis arcu lobortis. Aliquam in enim semper, aliquam massa id, cursus neque. Praesent faucibus semper libero.

% % Place figure captions after the first paragraph in which they are cited.
% \begin{figure}[!h]
% \caption{{\bf Bold the figure title.}
% Figure caption text here, please use this space for the figure panel descriptions instead of using subfigure commands. A: Lorem ipsum dolor sit amet. B: Consectetur adipiscing elit.}
% \label{fig1}
% \end{figure}

% Either type in your references using
% \begin{thebibliography}{}
% \bibitem{}
% Text
% \end{thebibliography}
%
% or
%
% Compile your BiBTeX database using our plos2015.bst
% style file and paste the contents of your .bbl file
% here.
% 
% \begin{thebibliography}{10}

% % \bibitem{bib1}
% % Conant GC, Wolfe KH.
% % \newblock {{T}urning a hobby into a job: how duplicated genes find new
% %   functions}.
% % \newblock Nat Rev Genet. 2008 Dec;9(12):938--950.

% % \bibitem{bib2}
% % Ohno S.
% % \newblock Evolution by gene duplication.
% % \newblock London: George Alien \& Unwin Ltd. Berlin, Heidelberg and New York:
% %   Springer-Verlag.; 1970.

% % \bibitem{bib3}
% % Magwire MM, Bayer F, Webster CL, Cao C, Jiggins FM.
% % \newblock {{S}uccessive increases in the resistance of {D}rosophila to viral
% %   infection through a transposon insertion followed by a {D}uplication}.
% % \newblock PLoS Genet. 2011 Oct;7(10):e1002337.
% \documentclass[12pt, final]{book}
%\documentclass[12pt, final]{book}
%\documentclass[12pt, draft]{book}
%\usepackage[index, nonumber]{cuisine}
%\usepackage[loose]{units}
\usepackage{url}
%\usepackage{path}
\usepackage{graphicx}
\usepackage{layout}
\usepackage{caption}
%\RecipeWidths{\textwidth}{0.01\textwidth}{0.2\textwidth}{0.25\textwidth}{0.09\textwidth}{0.08\textwidth}
% total, recipe #, servings, ingredients, quantity, units
%\renewcommand*{\recipetitlefont}{\large\bfseries\sffamily}
%\renewcommand*{\recipetex\titlefont}{\large\bfseries\sffamily}
%\renewcommand*{\recipequantityfont}{\sffamily\bfseries}
%\renewcommand*{\recipequantityfont}{\sffamily}
%\renewcommand*{\recipeunitfont}{\sffamily}
%\renewcommand*{\recipeingredientfont}{\sffamily}
%\renewcommand*{\recipefreeformfont}{\itshape}
\newcommand{\jalapeno}{jalape\~{n}o}
\newcommand{\jalapenos}{jalape\~{n}os}
\newcommand{\apf}{all purpose flour}
\newcommand{\bp}{baking powder}
\newcommand{\bs}{baking soda}
\newcommand{\ws}{Worcestershire sauce}
\newcommand{\tm}{\textsuperscript{\texttrademark}}

\newcommand{\beq}{\begin{equation}}
\newcommand{\eeq}{\end{equation}}


\newcommand{\metWalk}{\hat{\mathbf{\mathrm{MET}}}_{w}}
\newcommand{\metCycle}{\hat{\mathbf{\mathrm{MET}}}_{c}}
\newcommand{\metTransport}{\hat{\mathbf{\mathrm{MET}}}_{\mathrm{transport}}}
\newcommand{\relativeWalk}{\hat{\mathbf{R}}^t_w}
\newcommand{\relativeCycle}{\hat{\mathbf{R}}^t_c}
\newcommand{\speedCycle}{\hat{\mathbf{R}}^s_c}
\newcommand{\speedWalk}{\hat{\mathbf{R}}^s_w}

\newcommand{\popDensity}{\hat{\mathbf{F}}}
\newcommand{\sijw}{{\hat{s}_w}}

\newcommand{\walkMean}{\mathbf{T}_w}
\newcommand{\cycleMean}{\mathbf{T}_c}

\newcommand{\walkSpeedMean}{\mathbf{S}_w}
\newcommand{\cycleSpeedMean}{\mathbf{S}_c}

\newcommand{\distWalkMean}{\mathbf{D}_w}
\newcommand{\distCycleMean}{\mathbf{D}_c}
\newcommand{\RR}{\textrm{R.R.}}
\newcommand{\RRMET}{\RR_{\scriptsize{\textrm{MET}}}}
\newcommand{\RRMETi}{\RR_i^{\scriptsize{\textrm{MET}}}}
\newcommand{\RRvec}{\mathbf{\varphi}}
\newcommand{\burdenVec}{\mathbf{\xi}}
\newcommand{\quintMat}{\mathbf{Q}}

\begin{document}
%\layout
\frontmatter
\begin{titlepage}
\vspace*{1.5in}
\begin{center}
\resizebox{\linewidth}{!}{\textbf{ITHIM}}\\
\vspace*{0.5in}
\LARGE{\slshape Samuel Younkin. Ph.D.}\\
\vspace*{0.5in}
\Large{\itshape May, 2016}
%\vspace*{1in}
%\subtitle{
%\author{Sam \& Julie Younkin}
%\date{2013}
%\maketitle

% Upper part of the page
%\textsc{\LARGE University of Beer}\\[1.5cm]
%\textsc{\Large Final year project}\\[0.5cm]
% Title
%{ \huge \bfseries Lager brewing techniques}\\[0.4cm]
% Author and supervisor
%% \begin{minipage}{0.4\textwidth}
%% \begin{flushleft} \large
%% \emph{Author:}\\
%% John \textsc{Smith}
%% \end{flushleft}
%% \end{minipage}
%% \begin{minipage}{0.4\textwidth}
%% \begin{flushright} \large
%% \emph{Supervisor:} \\
%% Dr./home/sgy/Mark \textsc{Brown}
%% \end{flushright}
%% \end{minipage}

%% \vfill

%% % Bottom of the page
%% {\large \today}

\end{center}

\end{titlepage}

\cleardoublepage
\section*{Introduction}
Hello world.

\tableofcontents
\mainmatter

\chapter{Model}\label{ithimModel}

\section{Input Parameters}

\subsection{Population}
\beq
\popDensity = \begin{pmatrix}{f_{ij}}\end{pmatrix},\quad i \in \{\mathrm{M},\mathrm{F}\}, j \in \{\mathrm{age\ classes}\}  
\eeq

$\mu^t_w$ is the population mean time walking, with $\mu^t_c$, $\mu^s_w$, $\mu^s_c$ defined similarly with speed ($s$), time ($t$), walking ($w$) and cycling ($c$).

\subsection{Relative Time and Speed for Walking and Cycling}

\beq
\relativeWalk = \begin{pmatrix}{{\rho^t_w}_{ij}}\end{pmatrix}
\eeq

as well as $\relativeCycle$, $\speedWalk$ and $\speedCycle$ defined likewise.

\section{Model}

See \textit{Baseline} worksheet.

\subsection{Mean Walking and Cycling Times (minutes)}

\begin{eqnarray}
%%\hat{\mathbf{M}}_w &=& \begin{pmatrix}{\mu_{ij}}\end{pmatrix}\\
%%\mu_{ij} &=& \frac{\mu_w^{0}}{\alpha}{\rho_w}_{ij}\\
\walkMean &=& \frac{\mu^t_w}{\alpha^t_w}\relativeWalk, \quad \alpha^t_w = \sum_{ij} f_{ij}{\rho^t_w}_{ij}\\
\cycleMean &=& \frac{\mu^t_c}{\alpha^t_c}\relativeCycle, \quad \alpha^t_c = \sum_{ij} f_{ij}{\rho^t_c}_{ij}
\end{eqnarray}

\subsection{Mean Walking and Cycling Speeds}

\begin{eqnarray}
%%\hat{\mathbf{M}}_w &=& \begin{pmatrix}{\mu_{ij}}\end{pmatrix}\\
%%\mu_{ij} &=& \frac{\mu_w^{0}}{\alpha}{\rho_w}_{ij}\\
\walkSpeedMean &=& \frac{\mu^s_w}{\alpha^s_w}\speedWalk, \quad \alpha^s_w = \sum_{ij} f_{ij}{\rho^s_w}_{ij}\\
\cycleSpeedMean &=& \frac{\mu^s_c}{\alpha^s_c}\speedCycle, \quad \alpha^s_c = \sum_{ij} f_{ij}{\rho^s_c}_{ij}
\end{eqnarray}


\subsection{METs}
\beq
\metTransport = \metWalk + \metCycle
\eeq

\beq
\metWalk = f_w\left(\speedWalk\right), \quad f_w\left( \sijw \right) = \max\{ \beta_0 + \beta_1 \sijw, 2.5 \}
\eeq

with $\beta_1 = 1.2216$, $\beta_0 = 0.0838$ 

\beq
\metCycle = \mathbf{6}
\eeq

\subsection{\textit{Relative} Mean Distance of Active Transport}

Not sure about the relative distance measure.  What are the units?  Is
it used later for road injury miodeling.  I don't think it is used
here.

\beq
\distWalkMean = \frac{1}{60} \speedWalk \cdot \walkMean
\eeq
\beq
\distCycleMean = \frac{1}{60} \speedCycle \cdot \cycleMean
\eeq

\subsection{Quintiles}

Estimate quintiles per group based on group mean $\mu_{ij}$ and
constant coefficient of variation, $\gamma$ using a with $T$
distributed as lognormal.  Divided the quintiles proportionally into
walking and cycling, based on within-group proportions.  Compute MET-hours per quintile as $\metWalk T_w$

And we arrive at quintiles for total MET-hours per week stratified by age group and sex, $\quintMat$.



\section{Relative Risk given two types of Exposure (Weekly Active Time and MET-hours)}

%% \begin{verbatim}
%% =IF(($'user page'.$N$39=0),C5^(1/C6^0.25),
%% IF(($'user page'.$N$39=1),C5^(1/C6^0.5),
%% IF(($'user page'.$N$39=2),C5^(1/C6^0.375),IF(($'user page'.$N$39=4),C5^(1/C6),
%% IF(($'user page'.$N$39=3),C5^(1/(LN(C6))),""
%% )))))
%% \end{verbatim}

Let $x$ be the ``exposure'' in \textit{Phy Activity RRs} and $\RR$ the ``RR'' in the same worksheet then the ``1 MET'' described in the worksheet is named $\RR_{\scriptsize{\textrm{MET}}}$ here.

$$\RRMET = {\RR}^{{x^{-k}}}$$ with $k$ a model parameter.

\subsection{Definitions of $\RR$, Exposure ($x$) and $\RRMET$}

Do not have these yet.

For developmental purposes we fix the vector $\RRvec$ to be
the vector of disease sprecific relative risks given an exposure of 1
MET-hour per week.

$\RRvec = \left( \RRMETi \right)$

\subsection{Disease Burden}

First we define a weight vector for disease burden, $\burdenVec$. With
$\RRvec$, $\quintMat$ and $\burdenVec$ we can compute measures of
total disease burden.

\subsubsection{\textit{GBDUS}}

What are these data?  It's a table look up as with Non-travel METs.




\end{document}

%% \input{./../recipes/macaroni-and-cheese}
%% \newpage
%% \input{./../recipes/caesar-salad}
%% \newpage
%% \input{./../recipes/herb-pasta}
%% \newpage
%% \input{./../recipes/prosciutto-wrap}
%% \newpage
%% \input{./../recipes/lobster-grilled}
%% \newpage
%% \input{./../recipes/lobster-stock}
%% \newpage
%% \section{Breads}
%% \input{./../recipes/buttermilk-biscuits}
%% \newpage
%% \input{./../recipes/banana-bread}
%% \newpage
%% \input{./../recipes/monkey-bread}
%% \newpage
%% \section{Chili}
%% \input{./../recipes/beef-chili}
%% \newpage
%% \input{./../recipes/prison-chili}
%% \newpage
%% \input{./../recipes/cincinnati-chili}
%% \newpage
%% \section{Barbecue}
%% \input{./../recipes/basic-rib-rub}
%% \newpage
%% \input{./../recipes/barbecue-mop}
%% \newpage
%% \section{Desserts}
%% \input{./../recipes/allendale-brownies}
%% \newpage
%% \input{./../recipes/caramel-peach-crunch}
%% \newpage
%% \input{./../recipes/peanut-blossoms}
%% \newpage
%% \input{./../recipes/hough-frosting}
%% %% \newpage
%% %% \section{Coffee}
%% %% \input{./../recipes/chemex}
%% %% \newpage
%% %% \input{./../recipes/hario-v60}
%% \clearpage
%% \chapter{Thanksgiving}\label{chapter2}
%% \begin{figure}[h]
%% \begin{center}
%% \includegraphics[width=0.6\textwidth]{/home/sgy/Dropbox/cookbook/figures/turkey}
%% \end{center}
%% \caption*{Turkey on the barbecue. Baltimore, 201?}
%% \end{figure}
%% \newpage
%% \input{./../recipes/sailors-duff}
%% \input{./../recipes/duff-sauce}
%% \newpage
%% \input{./../recipes/stuffing}
%% \newpage
%% \input{./../recipes/onion-casserole}
%% \newpage
%% \input{./../recipes/green-bean-casserole}
%% \newpage
%% \input{./../recipes/cornbread}
%% \newpage
%% \input{./../recipes/turkey-carcass-soup}
%% \newpage
%% \input{./../recipes/pecan-tassies}
%% \newpage
%% \chapter{Canning \& Sauerkraut}
%% \input{./../recipes/red-pepper-relish}
%% \newpage
%% \input{./../recipes/sauerkraut}
%% \newpage
%% \chapter{The Less Favorites}\label{chapter3}
%% \input{./../recipes/barbecue-baked-beans}
%% \newpage
%% \input{./../recipes/grandmas-rice}
%% \newpage
%% \input{./../recipes/chicken-chili}
%% \newpage
%% \input{./../recipes/joc-chili}
%% \newpage
%% \input{./../recipes/ginger-soy-sauce}
%% \newpage
%% \input{./../recipes/refried-beans}
%% \newpage
%% \input{./../recipes/french-toast}
%% %~~~~~~~~~~~~~~~~~~~~~~~~~~~~~~~~~~~~~~
%% \chapter{Chili Powder}
%% \input{./../recipes/chili-powder-10}
%% \input{./../recipes/chili-powder-9}
%% \input{./../recipes/chili-powder-8}
%% \input{./../recipes/chili-powder-7}
%% \input{./../recipes/chili-powder-6}
%% \input{./../recipes/chili-powder-5}
%% \input{./../recipes/chili-powder-4}
%% \input{./../recipes/chili-powder}
%% \begin{figure}
%% \begin{center}
%% \includegraphics[height=0.5\textwidth]{/home/sgy/Dropbox/cookbook/figures/datil-bite-me}
%% \end{center}
%% \caption*{Datil chili peppers on the plant at Joe's ``Bite Me Farms'' in St.\ Augustine, FL.}
%% \end{figure}
%% \begin{figure}
%% \begin{center}
%% \includegraphics[height=0.3\textwidth]{/home/sgy/Dropbox/cookbook/figures/datil-sgy-2}
%% \includegraphics[height=0.3\textwidth]{/home/sgy/Dropbox/cookbook/figures/datil-sgy}
%% \end{center}
%% \caption*{Datil chili peppers and whole cumin seed}
%% \end{figure}

% \end{thebibliography}
